%%%%%%%%%%%%%%%%%%%%%%%%%%%%%%%%%%%%%%%%%%%%%%%%%%%%%%%%%%%%%%
%%%%%%%%% MA-0461 Algebra Lineal II %%%%%%%%%%%%%%%%
%%%%%%%%%%%%%%%%%%%%%%%%%%%%%%%%%%%%%%%%%%%%%%%%%%%%%%%%%%%%%%
\newpage
\subsection*{MA-0461 Algebra Lineal II}
\addcontentsline{toc}{subsection}{MA-0461 Algebra Lineal II}

\hypertarget{MA0461}{}
\begin{refsection}
\begin{table}[H]
\centering{
\begin{tabular}{|ll|}
\hline
%\multicolumn{2}{|c|}{\textbf{Nivel de virtualidad del curso:} Virtual}\\ \hline
\textbf{Año:} I  & \textbf{Requisitos:} Ninguno \\ 
\textbf{Ciclo:} I  & \textbf{Correquisitos:} Ninguno \\ 
\textbf{Tipo de curso:} Teórico & \textbf{Horas presenciales por semana:}   \\ 
\textbf{Créditos:} 3 & \textbf{Horas de trabajo independiente por semana:}  \\ \hline
\end{tabular}
}
\end{table}



\subsubsection*{Descripción del curso}

Este es un curso del tercer ciclo de la carrera de Matemática. Constituye el segundo de una secuencia de dos cursos de álgebra lineal. 

Este es el segundo de dos cursos b\'asicos en \'algebra lineal. El \'algebra lineal es,
en su origen, el estudio de los sistemas de ecuaciones lineales. Sin embargo, resulta
sorpendente que ese objetivo, aparentemente simple, lleva a una enorme riqueza de
teoremas y lenguage que ve aplicaciones en todas las \'areas de la matem\'atica y de la
ciencia. Conceptos como matrices, espacios vectoriales y las transformaciones lineales
entre estos espacios surgen naturalmente del estudio de sistemas lineales y resultan ser
la forma apropiada de tratar mir\'iadas de problemas en diversas \'areas.

Este curso presupone un conocimiento de la teor\'ia b\'asica de espacios vectoriales
de dimensi\'on finita y las transformaciones lineales entre estos espacios, as\'i como herramientas
computacionales elementales como el \'algebra matricial y determinantes. En
el curso agregaremos una estructura adicional para los espacios vectoriales, a saber la
de producto escalar. Esto nos permitir\'a una mayor riqueza de an\'alisis y facilitar\'a la
interpretaci\'on geom\'etrica de los resultados. 

En particular, estudiaremos una clase especial de bases
de espacios vectoriales donde los vectores son ortogonales entre s\'i. Las consecuencias
de esta inocente propiedad van a ser enormes.
M\'as adelante en el curso estudiaremos la forma en que las transformaciones lineales act\'uan sobre
subespacios vectoriales. Esto se ve reflejado de la forma m\'as elegante en la teor\'ia
espectral para matrices sim\'etricas, pero tambi\'en generalizaremos ese resultado para
lograr una clasificaci\'on completa con la forma normal de Jordan. Las consecuencias
pr\'acticas, en matem\'aticas y cualquier otra parte, de estos resultados no pueden exagerarse, son ubicuas y en muchas ocasiones son la pieza fundamental.

\subsubsection*{Objetivos}

\textbf{General:} Aplicar las nociones, intuiciones y conceptos del álgebra lineal para la resolución de problemas. \\

\textbf{Específicos:} Al finalizar el curso se espera que la persona estudiante sea capaz de:

\begin{enumerate}
\item Demostrar resultados sobre espacios vectoriales, transformaciones lineales, polinomios, formas canónicas, y formas multilineales.
%\item Aplicar los conceptos básicos del álgebra lineal en la demostración de resultados y la resolución de problemas teóricos y prácticos.
\item Realizar cálculos con los distintos objetos del álgebra lineal. (SH01)
\item Resolver problemas a través de los conceptos abstractos y formales del álgebra lineal. (SH05)
\item Relacionar los conceptos de valores y vectores propios con la acción de un operador lineal sobre un espacio vectorial.
\item Experimentar con los distintos objetos del álgebra lineal a través de software computacional. (SH03)
\item Indagar en la literatura sobre temas que permitan complementar su formación.
\item Comunicar ideas matemáticas de manera clara y efectiva en forma oral y escrita.
\end{enumerate}

\subsubsection*{Contenidos}

\begin{enumerate}

\item \textbf{Polinomios (2 semanas):} 
\begin{enumerate}
    \item Definiciones, funciones polinomiales y espacios vectoriales de polinomios.
    \item Propiedades básicas: unicidad de coeficientes, algoritmo de la división, propiedades del grado.
    \item Ceros de polinomios, el teorema del factor.
    \item El Teorema Fundamental del Álgebra y factorización de polinomios.
\end{enumerate}

\item \textbf{Dualidad (1.5 semanas):} 
\begin{enumerate}
    \item Funcionales lineales y dualidad. El doble dual.
    \item La transpuesta de una transformación lineal.
\end{enumerate}

\item \textbf{Valores y vectores propios (2.5 semanas):} 
\begin{enumerate}
    \item Productos, sumas directas y cocientes de espacios vectoriales.
    \item Valores y vectores propios.
    \item El polinomio característico.
    \item Similaridad, diagonalización y la multiplicidad de los valores propios.
    \item Subespacios invariantes.
\end{enumerate}

\item \textbf{Espacios con producto interno (3 semana):} 
\begin{enumerate}
    \item Productos internos y normas, la desigualdad de Cauchy-Schwarz.
    \item Ortonormalidad y el proceso de Gram-Schmidt.
    \item El Teorema de Representación de Riesz.
    \item Transformaciones lineales y sus representaciones en bases ortonormales.
    \item La adjunta de una transformación lineal.
    \item La identidad de Parseval y la desigualdad de Bessel.
    \item Operadores y matrices unitarias.
    \item Complementos y proyecciones ortogonales.
\end{enumerate}

\item \textbf{Triangularización, diagonalización y formas canónicas (4 semanas):} 
\begin{enumerate}
    \item El Teorema de Triangularización de Schur y el Teorema de Cayley-Hamilton.
    \item El polinomio mínimo.
    \item La forma canónica de Jordan.
        \item Operadores y matrices normales.
    \item El Teorema Espectral.
\end{enumerate}

\item \textbf{Formas bilineales y multilineales (2 semanas):} 
\begin{enumerate}
    \item Formas bilineales y formas cuadráticas. Polarización.
    \item Diagonalización de formas cuadráticas.
    \item Funciones multilineales y determinantes.
\end{enumerate}

\end{enumerate}

\subsubsection*{Metodología}

\noindent \textbf{Lineamientos metodológicos:} La persona docente a cargo de este curso debe respetar las siguientes pautas:
\begin{enumerate}
    \item La presentación de los contenidos debe priorizar la formalización de los conceptos del álgebra lineal. En particular en este curso se debe asumir que las personas estudiantes  ya manejan las intuiciones y las habilidades operacionales sobre los conceptos básicos del cálculo diferencial e integral y del álgebra lineal.
    \item En este curso se aplican las bases de lógica y las estrategias de demostración vistas en el curso MA-02XX Introducción a las Demostraciones.
    \item Las demostraciones deben presentarse tomando en cuenta que están dirigidas a una persona estudiante principiante. En particular, se debe priorizar dar un mayor nivel de detalle por encima de la brevedad y la estética, cuando esto ayude a mejorar la comprensión de los temas.
    \item Utilizar el recurso tecnológico para reforzar distintos conceptos estudiados en el curso por medio de la visualización, cálculo y experimentación. En particular, se podría asignar alguna tarea o proyecto que requiera utilizar un sistema algebraico computacional (CAS) para realizar algunas de estas tareas.
    \item Cuando la naturaleza del contenido lo permita, se deben utilizar representaciones visuales que apoyen la comprensión de los temas.
\end{enumerate}

\textbf{Sugerencias metodológicas:} Adicionalmente, se tienen las siguientes recomendaciones para la persona docente a cargo de este curso:
\begin{enumerate}
    \item Aprovechar momentos adecuados del curso para motivar el estudio por medio de reseñas acerca del desarrollo histórico de la matemática y su importancia en aplicaciones.
    \item En la evaluación, se sugiere utilizar estrategias que promuevan el trabajo constante de parte de las personas estudiantes a lo largo del ciclo lectivo. Por ejemplo, esto se puede hacer por medio de tareas o quizzes.
    \item Para fomentar el desarrollo de las habilidades investigativas en el estudiantado, se sugiere la implementación de pequeños proyectos de investigación. Por ejemplo, pueden ser proyectos de investigación bibliográfica o también proyectos cortos donde se desarrolle un tema por medio de ejercicios dirigidos.
    \item Dado que hoy en día la mayoría del trabajo en matemática, tanto a nivel académico como profesional, es de naturaleza colaborativa, se sugiere a la persona docente implementar estrategias que promuevan el trabajo en equipo.
\end{enumerate}




\nocite{Apo85}
\nocite{HK71}
\nocite{Axl15}
\nocite{Cur93}
%\nocite{Gro19}
\nocite{Tre17}
\nocite{MM18}
\nocite{GH17}




\printbibliography[heading=subbibliography,section=\the\value{refsection}]
\end{refsection}
